% Section 1
\section{Introduction}\label{sec:intro}

\subsection{The problem}

Ecological inference (EI) asks how groups behaved when only aggregates are
observed. A unit is a polling place, a precinct, a county or a riding. For
each unit we see the shares of its members in $R$ groups (how its electorate
voted in an earlier election, or its ethnic composition) and the shares of its
members in $C$ outcome categories (how it voted in a later election, or whether
it turned out). We want the latent $R\times C$ table of rates: the share of the
members of group $r$ who chose outcome $c$, unit by unit and, above all, for
the electorate as a whole. Transition rates between elections, turnout by
group, and racially polarised voting are the standard examples.

The data link groups and outcomes only through an accounting identity: in each
unit, the outcome shares are the group shares weighted by the unknown rates.
A unit has $RC$ unknown rates, but its margins supply only $C-1$ free
equations. Without further structure the rates are not point-identified, and
neither is their electorate-wide average.

The classical answers resolve this in one of two ways. The first assumes enough
to obtain a single number. Goodman's ecological regression assumes that the
rates are the same in every unit~\citep{goodman1953}. King's method assumes a
truncated bivariate normal distribution for the unit rates~\citep{king1997}, and
hierarchical R$\times$C models make related distributional
assumptions~\citep{rjkt2001,greinerquinn2009}. These assumptions do the
identifying work, and the data cannot check them. When they fail, for example
because group behaviour depends on the composition of the unit, the answer can
be confidently wrong. The second answer assumes nothing beyond the accounting
identity and reports the widest possible range, the Duncan--Davis ``method of
bounds''~\citep{duncan1953}. It is always valid, but often too wide to be
useful.

\subsection{The aim of the bounding approach}

Partial identification takes the middle road~\citep{manski2003,tamer2010,molinari2020}.
It states an assumption, derives the set of values of the target that are
consistent with the data and that assumption (the \emph{identified set}), and
reports that set. When the assumption is weak the set is wide, but it is
credible. When the assumption is strong the set is narrow, but it rests on more.
The researcher can see exactly what each assumption buys.

\citet{jkst2019} (JKST below) carried out this programme for $2\times2$ tables.
Their assumption is \emph{linear contextual effects} (LCE): each group's
expected rate may depend on the composition of the unit, but only linearly. LCE
is much weaker than Goodman's constancy, which is the special case with no
contextual effect. JKST show that under LCE the aggregate regression of the
outcome share on the group share is quadratic, that it identifies all but one
number, and that requiring the implied rates to be probabilities bounds that
number and hence the target.

This paper extends the programme in three directions and sets out what can be
learned under each assumption, with honest uncertainty:
\begin{enumerate}
\item \textbf{General R$\times$C tables.} The unidentified part grows from one
number to a $(C-1)R(R-1)/2$-dimensional set, and the bound becomes a linear
program (LP).
\item \textbf{A ladder of assumptions.} Units are often nested in districts
(precincts in ridings, counties in states). District baselines protect against
aggregation bias across districts, and shared contextual effects carry the
information of the most informative district to all of them. Sign beliefs from
monotone ecological inference~\citep{egghm2025} add further, cheap restrictions.
Each assumption is a rung, and rungs are reported side by side.
\item \textbf{Inference.} Every bound is computed from an estimated regression.
Each endpoint is the tightest of many estimated candidate limits, so the
plug-in bound is biased inward, sometimes badly. Valid confidence intervals must
treat each endpoint as an intersection bound~\citep{clr2013}.
\end{enumerate}

\subsection{Contributions}

The results are of four kinds, and we keep them apart throughout. A result
labelled \emph{proved} has a complete proof in the supplementary appendix. A
\emph{cited} result applies a published theorem whose conditions we verify. A
\emph{conjecture} is stated precisely, supported by checks or simulation, and
not proved. An \emph{algorithm} is a procedure whose validity is not proved;
simulation evidence does not promote it to a theorem.

\paragraph{Identification in R$\times$C.} Under LCE the regression of the outcome
shares on the group shares identifies the symmetric part $Q$ of the
contextual-effect matrices and nothing about the skew part $K$
(Lemma~\ref{lem:decomp}). We characterise the sharp identified set $\ThI$ for
$K$ through conditional support-function inequalities (Theorem~\ref{thm:sharp}),
define two larger sets $\ThX$ and $\ThO$, prove the nesting
$\ThI\subseteq\ThX\subseteq\ThO$ (Proposition~\ref{prop:nesting}), and show that
the outer set $\ThO$ is exactly what the regression function can deliver
(Proposition~\ref{prop:sharpreg}). $\ThO$ is a polytope, and bounds on any linear
target are an LP (Proposition~\ref{prop:lp}). For two groups the LP has a closed
form (Proposition~\ref{prop:tworows}). In $2\times2$ the construction reproduces
JKST's published bounds exactly (Proposition~\ref{prop:jkst}).

\paragraph{Districts and the rung ladder.} Grouping units adds nothing without a
model that links groups (Proposition~\ref{prop:D0}). With such a model, the
sharpness theorem carries over to any mean model that is affine in its unknowns
(Theorem~\ref{thm:affine}), and the rungs form a ladder of nested identified sets
(Theorem~\ref{thm:ladder}). We give a closed form for shared contextual effects
with two groups, which shows how one district with lopsided units and an extreme
baseline tightens the bound everywhere (Proposition~\ref{prop:D3}), and the exact
condition under which the pooled model is valid (Proposition~\ref{prop:D4}).

\paragraph{Signs, reinforcement and monotonicity.} Beliefs about the sign of the
between-group and within-group (contextual) associations, and the belief that the
two reinforce each other, are accounting identities away from data-computable
quantities (Lemma~\ref{lem:M1}). On realised tables they give LP bounds that are
sharp for continuous unit tables (Proposition~\ref{prop:N1}). Under the pooled
linear rung each becomes one linear row in $K$, the restricted bound is sharp for
the regression (Theorem~\ref{thm:R3}), and a face test decides exactly when a
belief tightens a target (Proposition~\ref{prop:R2}). These sign rows are
implemented in the accompanying package as an \emph{experimental} module: the
results are proved for the pooled model only, and sharpness and inference under
district rungs are open. A monotone nonparametric rung gives a valid outer LP
(Proposition~\ref{prop:N2}); its sharp version is a conjecture
(Conjecture~\ref{conj:N2}).

\paragraph{Inference.} Each endpoint is a minimum of finitely many linear
functions of the regression coefficients (Lemma~\ref{lem:minlin}). The plug-in
endpoint has a non-normal limit with negative mean when candidates tie
(Theorem~\ref{thm:limit}, Proposition~\ref{prop:triangular}). We prove asymptotic
coverage for a conservative interval and for projection and fixed-superset
intervals (Theorem~\ref{thm:ci}), give a variance-adaptive allowance for the
realised rate (Proposition~\ref{prop:bd}), tests of the model
(Proposition~\ref{prop:tests}), a finite-sample interval valid in every
asymptotic regime (Theorem~\ref{thm:finite}), and inference with added sign rows
and either/or restrictions (Theorem~\ref{thm:R10}). The CLR-type interval with
simulated candidates, the package default, is narrower in grouped designs, but
its validity is Conjecture~\ref{conj:clr}.

\paragraph{Software and evidence.} The R package \pkg{eiPartialID} implements the
rungs, the intervals and the tests, and prints for each fit which result covers
its intervals. A Monte Carlo study with 1000 replications per design, a benchmark
with known truth from New Zealand elections, and simulations of the sign module
show where each interval is informative and what happens when an assumption is
false.

\subsection{Related work}

The bounds of \citet{duncan1953} are the model-free baseline. \citet{crossmanski2002}
study ecological inference with an instrument that makes some rates constant
across units and obtain a point; LCE is weaker, and the bounds here lie between
their data-alone region and their point. The random-set formulation of sharp
identified sets follows \citet{bmm2011}. The inference builds on intersection
bounds~\citep{clr2013}, directionally differentiable functionals
\citep{shapiro1991,fangsantos2019}, and the distinction between covering the
parameter and covering the identified set~\citep{imbensmanski2004,stoye2009}.
The sign and reinforcement beliefs come from \citet{egghm2025} (EGGHM below);
their identities hold for any bounded outcome and are used here as accounting
identities. The view of the sources of uncertainty follows
\citet{aaiw2023} and \citet{aiz2014}.

\subsection{Plan}

Section~\ref{sec:setting} sets out the data, the targets, the model and the
assumptions, each with a plain-language reading. Section~\ref{sec:rxc} gives
the identification results for R$\times$C tables under pooled LCE.
Section~\ref{sec:districts} adds districts and the rung ladder.
Section~\ref{sec:signs} adds sign, reinforcement and monotonicity beliefs.
Section~\ref{sec:inference} turns the identified sets into confidence intervals
and tests. Section~\ref{sec:package} describes the R package,
Section~\ref{sec:evidence} the numerical evidence and a worked example, and
Section~\ref{sec:open} the open questions. The supplementary appendix contains
the proofs in one notation.

In the main text each formal statement is followed by a plain-language reading
and, for the key results, a sketch of the argument. A sketch is a guide to the
proof, not a substitute for it.
